\documentclass{article}
\usepackage[T1]{fontenc}
\usepackage{lmodern}
\usepackage{graphicx}
\usepackage{amsmath,amssymb}
\usepackage[a4paper,margin=2cm]{geometry}
\usepackage[absolute,overlay]{textpos}
\setlength{\TPHorizModule}{1mm}
\setlength{\TPVertModule}{1mm}
\usepackage[indonesian]{babel}
\usepackage{hyperref}
\setlength{\emergencystretch}{2em}
% unit: Halaman 17

\begin{document}

% === Halaman 17 (llm) ===

\begin{itemize}
    \item Contoh di dalam Caesar Cipher, penyerang memilih plainteks:
    
    A \\ $\rightarrow$ Cipherteks keluar: F. Artinya pergeseran mungkin 5.

    Untuk memastikan, penyerang pilih plainteks lain:
    
    B \\ $\rightarrow$ Cipherteks keluar: G. Polanya konsisten, maka kunci pergeseran dipastikan = 5.
\end{itemize}

\end{document}
